% Background and related work. Every claim here is about prior work; results of this paper are in results.tex.
\section{Background and related work}
\label{sec:background}

Single-cell and spatial transcriptomics produce one kind of object above all others: a sparse matrix of non-negative
integer counts, genes by cells, in which most entries are zero and most non-zero entries are small. The number of
cells per study has grown exponentially for a decade~\cite{svensson2018exponential}, reference atlases hold millions of
cells~\cite{tabulasapiens2022,sikkema2023integrated}, public collections approach a hundred
million~\cite{czi2025cellxgene}, and a single perturbation experiment has reached a hundred
million~\cite{zhang2025tahoe}. These matrices are read far more often than they are written: by analysis
packages, by atlases that serve them to thousands of users, and increasingly by deep learning models that stream
them as minibatches for many epochs~\cite{lopez2018deep,theodoris2023transfer,cui2024scgpt}. This section reviews how
such matrices are stored, how integer data are compressed in other fields, what is known about the statistics of the
counts, and how data reach a model. Together these fix what a storage format for count matrices can achieve and
where existing formats leave room.

\subsection{Sparse matrix storage}

The compressed sparse column (CSC) and row (CSR) layouts store, for each column, the row indices and values of its
non-zero entries, with one pointer per column marking where its entries begin~\cite{saad2003iterative,davis2006direct}.
They are the in-memory representation of almost every single-cell toolkit: \texttt{dgCMatrix} in the R
\texttt{Matrix} package~\cite{bates2024matrix}, SciPy sparse matrices~\cite{virtanen2020scipy}, and the sparse tensors
of PyTorch~\cite{paszke2019pytorch,ansel2024pytorch2}. On disk a CSC matrix is three arrays, typically 32-bit row indices and 32- or
64-bit values, so it costs 64 to 96 bits per non-zero before any compression, however small the counts are.

Exchange formats wrap these arrays. Matrix Market~\cite{boisvert1997matrix} writes coordinate triplets as text and is
what Cell Ranger and many GEO submissions distribute, usually gzip-compressed. Hierarchical containers store the
arrays as chunked, optionally compressed datasets: the 10x Genomics HDF5 layout of Cell
Ranger~\cite{zheng2017massively,tenx_h5}, loom (dense chunks in HDF5)~\cite{loompy}, and AnnData~\cite{virshup2024anndata}, whose \texttt{.h5ad} (HDF5~\cite{hdf5}) and Zarr~\cite{zarr}
encodings are the standard of the Python ecosystem built on Scanpy~\cite{wolf2018scanpy}. In R the same role is
played by serialized \texttt{dgCMatrix} objects inside SingleCellExperiment~\cite{amezquita2020orchestrating} and
Seurat~\cite{hao2024dictionary} objects. In every case compression is a generic byte-level filter (gzip/DEFLATE~\cite{deutsch1996deflate},
LZF, Blosc, or zstd~\cite{collet2021zstd}) applied to the arrays, which knows nothing of their structure.

Formats designed for scale add structure. TileDB-SOMA stores matrices as sparse arrays in tiles with per-attribute
filter pipelines~\cite{papadopoulos2016tiledb,tiledbsoma}, and serves the CZ CELLxGENE Census of tens of millions of
cells~\cite{czi2025cellxgene}; cells can be queried by metadata and read without loading the whole matrix.
BPCells~\cite{parks2025bpcells} stores row indices as bit-packed deltas and values as bit-packed integers, in the spirit
of SIMD integer codecs (below), and reads them through a streaming iterator, so operations need far less memory than
loading the matrix; it reports files smaller than \texttt{dgCMatrix} and fast streaming reads. Parquet
tables~\cite{parquet} are used to distribute the largest training corpora~\cite{tahoe100m_hf}. Of these, only BPCells
uses a transform tailored to count data, and none uses an entropy coder with a model of the counts.

\subsection{Compressing integers}

Lossless compression separates modelling, which assigns a probability to each symbol, from coding, which spends close
to $-\log_2 p$ bits on it~\cite{shannon1948mathematical,cover2006elements}. Huffman codes~\cite{huffman1952method} are
optimal among prefix codes but can lose up to a bit per symbol on skewed distributions~\cite{cover2006elements}, which
is exactly the regime of count data, where most values are 1. Arithmetic
coding~\cite{rissanen1979arithmetic,witten1987arithmetic} removes that loss at a much larger computational
cost~\cite{duda2013asymmetric}. Asymmetric numeral systems (ANS)~\cite{duda2009asymmetric,duda2013asymmetric} reach the
precision of arithmetic coding with table lookups, and the range variant rANS can run many independent states
interleaved in one stream~\cite{giesen2014interleaved}, so that SIMD units decode several symbols per instruction. A
tabled ANS variant codes all but the literals in zstd~\cite{collet2021zstd}, rANS with static tables is an entropy
coder of CRAM~3.1~\cite{bonfield2022cram}, and the format described here uses it too.

Inverted indexes in information retrieval face the same problem as CSC row indices: sorted integer lists whose gaps
are small. Classical codes are Elias $\gamma$ and $\delta$~\cite{elias1975universal}, Golomb and Rice
codes~\cite{golomb1966run,rice1979some}, and Elias--Fano~\cite{elias1974efficient,vigna2013quasi}. For speed, byte- and word-aligned codes
replaced bit-level ones: patched frame-of-reference (PFOR)~\cite{zukowski2006super}, SIMD bit packing that decodes
billions of integers per second~\cite{lemire2015decoding}, Stream VByte~\cite{lemire2018stream}, and hybrid
bitmap/array containers such as Roaring~\cite{chambi2016better}; partitioned Elias--Fano~\cite{ottaviano2014partitioned}
and binary interpolative coding~\cite{moffat2000binary} trade speed for size~\cite{pibiri2021techniques}. Most of these
codes are fast because they ignore the distribution of the values beyond their bit width, leaving the entropy of skewed
gap and value distributions on the table; ANS with contexts closes much of that gap for inverted
indexes~\cite{moffat2017ans,moffat2018index,moffat2020large}, the closest precedent for the approach taken here. Generic
filters such as Blosc with bit shuffling~\cite{alted2010modern,masui2015compression,alted2023blosc2} and compressors primed
with a trained dictionary~\cite{collet2016smaller,liao2016effective} help on small, similar blocks without modelling the
values. General-purpose compressors (DEFLATE, bzip2~\cite{burrows1994block}, LZMA/xz, brotli~\cite{alakuijala2016brotli},
zstd~\cite{collet2021zstd}) find repeated byte strings and model bytes with small contexts; they compress the
\emph{representation} of a sparse matrix rather than the counts. Context-mixing compressors such as
PAQ~\cite{mahoney2005adaptive} come closer to the entropy of almost any byte stream, but run hundreds to thousands of
times slower than gzip.

\subsection{Compression of genomic data}

Genomics has a long record of domain-specific formats that beat generic compression by modelling the data.
Reference-based compression of aligned reads~\cite{fritz2011efficient} led to CRAM~\cite{cochrane2013facing}, which
replaced BAM's gzip blocks with per-field codecs; CRAM~3.1 added rANS with optionally 32 interleaved states and SIMD
decoding, a context-model quality coder and a name tokeniser, reaching files 50--70\% smaller than BAM on Illumina
data~\cite{bonfield2022cram}. Specialised FASTQ compressors model
sequences, qualities and read names separately~\cite{bonfield2013compression,chandak2019spring}, and benchmarks across
dozens of tools found that how each field is modelled matters more to the ratio than the choice of back-end
coder~\cite{numanagic2016comparison,hernaez2019genomic}.
The lessons that carry over are: separate fields with different statistics into their own streams, condition each
on cheap side information, use static tables learned from the data being coded, and keep random access by coding in
independent blocks. Genotype matrices, sparse integer matrices of a different kind, have dedicated compressors that
reorder and model them~\cite{durbin2014efficient,danek2018gtc,deorowicz2019gtshark,lan2021genozip}, as reordering does
for graphs~\cite{boldi2004webgraph,boldi2011layered,dhulipala2016compressing}. Count matrices have received far less
attention: we found no published entropy-coded format for them, and the formats in use rely on generic filters.

\subsection{What the counts look like}

A UMI count $x_{gc}$ for gene $g$ in cell $c$ is well described as a sample from the cell's pool of molecules. If cell
$c$ contributes $N_c$ captured molecules with expression proportions $\pi_c$, the counts are multinomial,
$x_c \sim \mathrm{Multinomial}(N_c, \pi_c)$, which for a large pool is equivalent to independent Poisson counts given
$N_c\pi_c$~\cite{townes2019feature}. Technical noise in UMI data is close to this Poisson sampling
limit~\cite{kim2020demystifying,townes2019feature,sarkar2021separating}, with an additional cell-to-cell variation in
capture efficiency~\cite{grun2014validation}, and droplet data are not zero-inflated once biological variation is
modelled~\cite{svensson2020droplet,kim2020demystifying}: the excess of zeros over a Poisson is explained by variation of
$\pi$ across cells, which a gamma-Poisson (negative binomial) mixture captures~\cite{svensson2020droplet,choudhary2022comparison}.
The negative binomial with a gene-specific dispersion, long the workhorse of bulk
RNA-seq~\cite{robinson2010edger,love2014moderated}, underlies sctransform~\cite{hafemeister2019normalization} and the
likelihoods of deep generative models~\cite{lopez2018deep,eraslan2019single}, while GLM-PCA models counts as multinomial
or Poisson~\cite{townes2019feature}. Zero-inflated variants~\cite{risso2018general} matter for read-count protocols such
as Smart-seq2, where PCR duplication inflates counts~\cite{townes2019feature}.
Expression across genes is heavy-tailed and close to Zipf's law~\cite{furusawa2003zipf}, the same rank-frequency
law as word counts in text, which is why bag-of-words corpora are a natural out-of-domain test.

These models matter for compression in two ways. First, they say what an entropy coder can exploit: the depth of a cell
($N_c$) sets the rate of every gene in it; a gene's mean sets its presence probability; and most non-zero values are 1.
Second, they bound what any coder can achieve. Under the multinomial model the code length of a cell splits into
the entropy of sampling $N_c$ molecules from $\pi_c$, which no decoder can avoid because $\pi_c$ is the content being
transmitted, and the information the counts carry about $\pi_c$ itself~\cite{cover2006elements}. The second term is the
only one a better model, more side information, or another cell can reduce. The experiments in
Section~\ref{sec:sfloor} measure both terms.

\subsection{Getting counts into a model}

Training on single-cell data means drawing random minibatches of cells, many times over. Loading a whole matrix into
memory is simplest and is what most pipelines do, but it fails beyond the memory of one node. Backed AnnData reads
rows of an on-disk \texttt{.h5ad}, \texttt{AnnLoader} wraps AnnData for PyTorch, and atlas-scale projects built their own
loaders: scTab used NVIDIA Merlin over Parquet~\cite{fischer2024scaling}, the CELLxGENE Census provides
TileDB-SOMA-ML~\cite{tiledbsomaml,czi2025cellxgene}, and scDataset showed that fetching blocks of contiguous cells, then
shuffling within a larger buffer, is more than two orders of magnitude faster than fully random sampling, with the same
model accuracy on the tasks tested~\cite{dascenzo2025scdataset}. These loaders optimise I/O patterns over existing
formats. When the data sit on a network file system or in an object store, as in most clusters and clouds, two costs set
the ceiling on throughput: the bytes moved per cell, which the format sets, and the number of round trips, each of which
waits tens of milliseconds on a cloud object store~\cite{durner2023exploiting}. Cloud-native scientific formats such as
Zarr and OME-Zarr address the second by splitting data into independently addressable chunks that can be fetched
concurrently~\cite{abernathey2021cloud,moore2021omengff}.

Whether compression helps a loader depends on both the bytes per value $b$ and the decode rate $D$. Delivered over a link of
bandwidth $B$, a format provides values at the effective rate $R_{\mathrm{eff}} = (b/B + 1/D)^{-1}$ (Eq.~\ref{eq:reff}),
so a smaller format is only better if it decodes fast enough that $1/D$ does not eat its transfer savings. The model
is conservative: it assumes transfer and decoding happen one after the other, while a reader that overlaps them does
better. Formats with generic compression shrink $b$ modestly and decode at tens of millions of values per second in
practice; uncompressed formats have the highest $D$ and the largest $b$. The question this paper asks is whether a count
matrix can be both much smaller than existing formats and fast enough to decode that the network, not the decoder,
stays the limit on the links data actually travel over.
